% !TeX root = ../main.tex
\section{Appendix: Parameter Provenance}
\label{deprecated-sec:appendix-provenance}

This appendix consolidates parameter provenance and sensitivity checks. Proxy estimates are labelled explicitly.

\subsection{Three-Seed Summary for the Canonical V4 Path}
\label{deprecated-subsec:v4-rerun-sanity}

Table~\ref{deprecated-tab:v4-three-seed-summary} summarizes protocol-matched fresh-instance reproducibility across three Ensemble HAT training seeds under the canonical regime.

\begin{table}[ht]
\centering
\caption{Three-seed fresh-instance reproducibility summary for the canonical Ensemble HAT regime. Each seed is evaluated on 10 fresh D2D instances with five Monte Carlo passes per instance.}
\label{deprecated-tab:v4-three-seed-summary}
\begin{tabular}{lcc}
\toprule
\textbf{Seed / Aggregate} & \textbf{Fresh Accuracy} & \textbf{Statistic Type} \\
\midrule
123 & \(86.37\pm1.54\)\% & 10 fresh instances \(\times\) 5 MC \\
456 & \(86.12\pm0.72\)\% & 10 fresh instances \(\times\) 5 MC \\
789 & \(85.99\pm1.94\)\% & 10 fresh instances \(\times\) 5 MC \\
\midrule
Cross-seed aggregate & \(\mathbf{86.16\pm0.19}\)\% & mean \(\pm\) sample std of seed means \\
Pooled fresh instances & \(86.16\pm1.52\)\% & mean \(\pm\) sample std over 30 instance means \\
\bottomrule
\end{tabular}
\footnotesize Source: \texttt{r10a\_canonical\_ensemble\_hat\_3seed\_fresh\_eval.json}. The \(86.37\pm1.54\%\) value is retained only for the original plotted checkpoint; the training-seed reproducibility headline is \(86.16\pm0.19\%\).
\end{table}

\FloatBarrier

\begin{table}[ht]
\centering
\caption{Parameter provenance tracking matrix.}
\label{deprecated-tab:provenance}
\scriptsize
\setlength{\tabcolsep}{1.8pt}
\renewcommand{\arraystretch}{1.1}
\begin{tabular}{>{\raggedright\arraybackslash}p{2.1cm}>{\raggedright\arraybackslash}p{1.9cm}>{\raggedright\arraybackslash}p{2.2cm}>{\raggedright\arraybackslash}p{2.0cm}>{\raggedright\arraybackslash}p{5.0cm}}
\toprule
\textbf{Parameter} & \textbf{Canonical Organic Profile} & \textbf{Zhang 2025 OPECT Case Study} & \textbf{Task 34-36 Stress Tests} & \textbf{Provenance / Derivation Notes} \\
\midrule
 Conductance Range (\(G_{\max}/G_{\min}\)) & 10 & 47.3 & 10 & Canonical anchored to scalable OPECTs; Zhang derives from reported max/min current levels \citep{liu2026opect}. \\
 Effective States ($n_{states}$) & 16 (4-bit) & 34 & 16 (4-bit) & Canonical matches general low-precision target. Zhang anchored to Fig.3h \& Supp.Fig.8 \citep{liu2026opect}. \\
 Cycle-to-Cycle Noise (\(\sigma_{\text{C2C}}\)) & 5\% & 2\% & 5\% (Proportional) & Zhang value is a proxy estimate from 8-cycle repeatability (Supp.Fig.15) \citep{liu2026opect}. Proportional stress scales \(\sigma \propto |G|\). \\
  Device-to-Device Variability (\(\sigma_{\text{D2D}}\)) & 10\% & 3\% & 10\% (Proportional) & Zhang uses 3\% as a conservative conductance-domain proxy from the reported $\sim$1\% $V_{th}$ spread \citep{liu2026opect}. \\
 Retention Decay (\(\tau_1, \tau_2, A_0\)) & 140ms, 610ms, 0.6 & Not fitted & 140ms, 610ms, 0.6 & Canonical dual-exponential fit anchored to Vincze \citep{vincze2025dualplasticity}. Zhang (Fig.2d) is qualitative only. \\
 Nonlinear Write Asymmetry (\(NL\)) & 1.0 (LTP), -1.0 (LTD) & Canonical defaults retained & 2.0 (LTP), -2.0 (LTD) & Zhang does not report a pulse-level NL fit, so the case study keeps the canonical symmetric-write defaults rather than injecting a separate literature-fit NL term. Stress tests (\(NL=2\)) enforce severe gradient-scaling asymmetry. \\
 Noise Mode & Uniform & Uniform & Proportional & Uniform injects noise normalized to full range. Proportional applies state-dependent noise magnitudes. \\
\bottomrule
\end{tabular}
\end{table}

\FloatBarrier

\subsection{Vincze et al. 2025 Parameter Extraction}
\label{deprecated-subsec:vincze-extraction}



\subsection{Proxy Estimate Sensitivity Analysis}
\label{deprecated-subsec:proxy-sensitivity}

Table~\ref{deprecated-tab:sensitivity} validates the Zhang 2025 OPECT proxy estimates via 2D noise sweep.

\begin{table}[ht]
\centering
\caption{Ensemble HAT accuracy under varied Zhang proxy noise assumptions. The repeated rows across the C2C sweep are expected under the present checkpoint because the deployment accuracy is dominated by static D2D variability; the sampled C2C term contributes below the reported precision after Monte Carlo averaging.}
\label{deprecated-tab:sensitivity}
\small
\begin{tabular}{lccccc}
\toprule
\textbf{C2C \textbackslash D2D} & \textbf{2\%} & \textbf{3\% (Nominal)} & \textbf{5\%} & \textbf{10\%} & \textbf{15\%} \\
\midrule
\textbf{1\%} & 88.57\% & 88.53\% & 88.33\% & 87.30\% & 84.59\% \\
\textbf{2\% (Nominal)} & 88.57\% & 88.53\% & 88.33\% & 87.30\% & 84.59\% \\
\textbf{5\%} & 88.57\% & 88.53\% & 88.33\% & 87.30\% & 84.59\% \\
\textbf{8\%} & 88.57\% & 88.53\% & 88.33\% & 87.30\% & 84.59\% \\
\bottomrule
\end{tabular}
\end{table}

Static D2D spatial variability dominates the case-study outcome; per-forward C2C noise remains sub-resolution.

\FloatBarrier

\subsection{Uniform vs. State-Dependent Retention Comparison}
\label{deprecated-subsec:retention-comparison}

Table~\ref{deprecated-tab:retention-comparison} compares uniform and state-dependent retention decay.

\begin{table}[ht]
\centering
\caption{Retention accuracy comparison: uniform vs. state-dependent decay.}
\label{deprecated-tab:retention-comparison}
\begin{tabular}{lccc}
\toprule
\textbf{Time (s)} & \textbf{Uniform (\%)} & \textbf{State-Dep (\%)} & \textbf{Diff (pp)} \\
\midrule
0    & 90.77 & 90.80 & +0.03 \\
1    & 90.07 & 90.08 & +0.01 \\
10   & 89.78 & 89.81 & +0.03 \\
100  & 89.87 & 89.89 & +0.02 \\
1000 & 89.82 & 89.90 & +0.08 \\
\bottomrule
\end{tabular}
\end{table}

The accuracy difference remains below 0.1~pp, validating the uniform-retention assumption for the Ensemble-HAT regime.

\FloatBarrier

\subsection{Automated Profile Fitting Pipeline}
\label{deprecated-subsec:auto-fitting}

Device parameters can be extracted from measured programming, variability, and retention data via regression (Supplementary Methods).
